\textbf{R:} Para calcular los valores de acción, partiremos de lo que ya sabemos y simplificaremos un poco la fórmula, podemos asumir que como todos los valores $Q$ iniciales son $0$, eso significa que la máxima recompensa futura será $0$ para cualquier $s'$.

Tenemos entonces:
\begin{align*}
    Q(s, a) &= Q(s, a) + \alpha [r(s,a) + \gamma \max_{a'} Q(s', a') - Q(s, a)] \\
    &= 0 + \alpha [r(s,a) + \gamma(0) - 0] \\
    &= \alpha \cdot r(s,a)
\end{align*}

Ahora veamos las transiciones posibles:
\begin{itemize}
    \item Estado $0$
          \begin{itemize}
              \item $Q(0, Derecha) = \alpha \cdot r = \alpha \cdot 0 = 0$
              \item $Q(0, Abajo) = \alpha \cdot r = \alpha \cdot 0 = 0$
          \end{itemize}

    \item Estado $1$
          \begin{itemize}
                \item $Q(1, Izquierda) = \alpha \cdot r = \alpha \cdot 0 = 0$
                \item $Q(1, Derecha) = \alpha \cdot r = \alpha \cdot 0 = 0$
                \item $Q(1, Abajo) = \alpha \cdot r = \alpha \cdot 0 = 0$  
          \end{itemize}
    \item Estado $2$
            \begin{itemize}
                \item $Q(2, Izquierda) = \alpha \cdot r = \alpha \cdot 0 = 0$
                \item $Q(2, Abajo) = \alpha \cdot r = \alpha \cdot 1 = \alpha$
            \end{itemize}
    \item Estado $3$
            \begin{itemize}
                \item $Q(3, Arriba) = \alpha \cdot r = \alpha \cdot 0 = 0$
                \item $Q(3, Derecha) = \alpha \cdot r = \alpha \cdot 0 = 0$
            \end{itemize}
    \item Estado $4$
            \begin{itemize}
                \item $Q(4, Izquierda) = \alpha \cdot r = \alpha \cdot 0 = 0$
                \item $Q(4, Derecha) = \alpha \cdot r = \alpha \cdot 1 = \alpha$
                \item $Q(4, Arriba) = \alpha \cdot r = \alpha \cdot 0 = 0$
            \end{itemize}
\end{itemize}

Como podemos apreciar en las trancisiones, para cualquier acción en los estados $0,1,3$ el valor de $Q$ es $0$, mientras que para las acciones $Abajo$ en el estado $2$ y $Derecha$ en el estado $4$, el valor de $Q$ es $\alpha$ (las demas acciones en estos estados tendran un valor para $Q$ de $0$).

\begin{comment}
    $Q(s, a) = Q(s, a) + \alpha [r(s,a) + \gamma \max_{a'} Q(s', a') - Q(s, a)]$
    r(s,a) = recompensa obtenida en ese estado dad auna solución.
    $\gamma$ = factor de descuento para un estado futuro, entre 0 y 1.
    $\alpha$ = tasa de aprendizaje, entre 1 e irá descendieondo hasta 0.
    $Q(s,a)$ = valor Q para el estado $s$ y la acción $a$.
    $max_{a'} Q(s', a')$ = máxima futura recompensa esperada para el estado $s'$.
\end{comment}